% opus_03_theoremA_K2_id.tex
% Wave-17 OPUS adversarial-heal record for Theorem A
%   (Omega_X, \bar B_X) : Coalg_X^{conil,comp} <-> Alg_X^{aug,comp}
%   symmetric-monoidal adjoint equivalence, with K^2 simeq id on Kosz(X).
%
% Role: adversarial audit + heal record for the unified Theorem A
% inscribed at chapters/theory/theorem_A_infinity_2.tex:158--244
% (thm:koszul-reflection). Voices: BEILINSON (falsification) +
% DRINFELD (unifying principle), with WITTEN for physical cross-checks.
%
% This file is a NOTE (notes/wave17_opus_20260424/), not a chapter.
% It records the 5+ attack-heal cycles and the distilled HEAL text
% to be merged back into theorem_A_infinity_2.tex in a subsequent pass.

\documentclass[11pt]{article}
\usepackage[localtheorems]{raeez-math-template}

\newtheorem{theorem}{Theorem}
\newtheorem{lemma}[theorem]{Lemma}
\newtheorem{proposition}[theorem]{Proposition}
\newtheorem{corollary}[theorem]{Corollary}
\theoremstyle{definition}
\newtheorem{definition}[theorem]{Definition}
\newtheorem{remark}[theorem]{Remark}

\newcommand{\Bbarch}{\bar B^{\mathrm{ch}}}
\newcommand{\Omegach}{\Omega^{\mathrm{ch}}}
\newcommand{\Kosz}{\mathsf{Kosz}}
\newcommand{\Fact}{\mathsf{Fact}}
\newcommand{\Conil}{\mathsf{Conil}}
\newcommand{\Ran}{\mathrm{Ran}}
\newcommand{\Tc}{T^c}
\newcommand{\id}{\mathrm{id}}
\newcommand{\Ho}{\mathrm{Ho}}
\newcommand{\cA}{\mathcal{A}}
\newcommand{\cC}{\mathcal{C}}
\newcommand{\cP}{\mathcal{P}}
\newcommand{\cV}{\mathcal{V}}
\newcommand{\fg}{\mathfrak{g}}
\newcommand{\Eone}{E_1}
\newcommand{\RHom}{\mathbf{R}\mathrm{Hom}}
\newcommand{\Tor}{\mathrm{Tor}}
\newcommand{\Ext}{\mathrm{Ext}}
\newcommand{\Hom}{\mathrm{Hom}}
\newcommand{\Ainf}{A_\infty}
\newcommand{\MC}{\mathrm{MC}}
\newcommand{\simequiv}{\xrightarrow{\sim}}

\title{Theorem A, audited: is $K^{2}\simeq\id$ on $\Kosz(C)$ a
  chain-level theorem, a Quillen equivalence, an $(\infty,1)$-adjoint
  equivalence, or an $(\infty,2)$-equivalence? \\
  \small Wave-17 adversarial-heal record (Opus, 2026-04-24)}
\author{Beilinson-voice attack + Drinfeld-voice heal}
\date{2026-04-24}

\begin{document}
\maketitle

\begin{abstract}
The Wave-14 draft at \texttt{chapters/theory/theorem\_A\_infinity\_2.tex}
states Theorem A as a single Koszul reflection
\[
K:\Alg_X^{\fact,\aug,\comp}\rightleftarrows\mathrm{CoAlg}_X^{\fact,\conil,\comp}:K^{-1},
\qquad K(\cA)=\Bbarch_X(\cA),
\qquad K^{-1}(\cC)=\Omegach_X(\cC),
\]
with $K^{2}\simeq\id$ on $\Kosz(X)$ and coderived involutivity off
$\Kosz(X)$. This note runs seven attack-heal cycles (A1--A7) and
records the healed text to be merged into the chapter. Chief
outcomes: (a) the chain-level / $(\infty,1)$ / $(\infty,2)$ lanes
carry three genuinely different $K^{2}\simeq\id$ statements (not one
theorem in three dresses); (b) the Francis--Gaitsgory properad-level
$(\infty,2)$-arrow does NOT by itself entail the chain-level Quillen
equivalence -- the chain-level upgrade requires Mittag-Leffler
(H3-supplied) plus $E_2$-collapse (PBW + Koszulness); (c) $\Kosz(X)$
is defined antecedently in \texttt{chiral\_koszul\_pairs.tex:957}
(Definition of chiral Koszul pair), so the statement is
non-circular; (d) on the four archetypes $\mathsf{G}/\mathsf{L}/\mathsf{C}/\mathsf{M}$
each representative satisfies $K^{2}\simeq\id$ with a named chain
homotopy at the bar level; (e) the MC4-completion regime supplies a
specific $\cA$ (critical-level affine Kac--Moody with unbounded
depth) where (H3) fails, so the chain-level $K^{2}\simeq\id$
DEGRADES to the coderived statement but does not fail -- the
coderived equivalence of \ref{KR-iii} is the honest invariant.
\end{abstract}

%================================================================
\section*{Attack-heal cycle log (seven cycles)}
%================================================================

\subsection*{Cycle 1 (BEILINSON). Chain-level vs $(\infty,1)$-homotopy: which lane does $K^{2}\simeq\id$ live in?}

\textbf{Attack.} The wave-14 chapter states $K^{2}\simeq\id$ on $\Kosz(X)$ in
clause \ref{KR-ii} with language ``chain-level quasi-isomorphisms for
both unit and counit'' (line 188--191). But the proof at line 203--244
routes through (a) Vallette bar--cobar Quillen equivalence
(\cite{Val16}), (b) Positselski covariantly-weighted coacyclicity
(\cite{Positselski2011}), (c) the seven-fold TFAE at
$E_2$-collapse hub (\texttt{ftm\_seven\_fold\_tfae\_platonic.tex}),
and (d) promotion from coderived to chain-level via H3. Each of
these lives in a different lane:
\begin{itemize}
\item Vallette is a Quillen equivalence between model categories
  (a HOMOTOPY-categorical equivalence, not a chain-level isomorphism).
\item Positselski is a comparison of coderived categories
  $D^{\co}\simeq D^{\ctr}$ (a triangulated-categorical statement).
\item The seven-fold TFAE is a genus-zero statement on $\Kosz(X)$.
\item The promotion to chain-level requires H3 finite-dim bar pieces.
\end{itemize}
Is the object labelled ``$K^{2}\simeq\id$ on $\Kosz(X)$'' a chain-level
theorem, a Quillen equivalence, an $(\infty,1)$-adjoint equivalence,
or an $(\infty,2)$-equivalence at properad level?

\textbf{Heal (DRINFELD).} The statement is FOUR genuinely different
theorems, all true, all labelled by the same schematic slogan:
\begin{enumerate}[label=\textup{(L\arabic*)}]
\item \label{L1} \emph{Chain-level.} On $\Kosz(X)$ satisfying (H1)+(H2)+(H3),
the unit $\eta_\cA\colon\cA\to\Omegach_X\Bbarch_X(\cA)$ and counit
$\varepsilon_\cA$ are chain-level quasi-isomorphisms, and each
archetype carries an explicit chain homotopy
$h_\cA\colon\cA\to\Omegach_X\Bbarch_X(\cA)[-1]$ with
$[d,h_\cA]=\id-p_\cA$ for the retraction
$p_\cA=\varepsilon_\cA\circ\eta_\cA^{-1}$ restricted to the stratum
where the retract splits (Heisenberg $\mathsf{G}$: $h$ is the
Koszul contraction along partition-number grading;
$V_k(\mathfrak{sl}_2)$ $\mathsf{L}$ at noncritical level: $h$ is the
Sugawara-twisted deformation of the Heisenberg contraction; Virasoro
$\mathsf{M}$: $h$ is the Felder BRST reduction chain homotopy on the
weight-completed $\widehat{\mathrm{Vir}}_c$; $\beta\gamma$ $\mathsf{C}$:
$h$ is the free-field OPE contraction along the $b\gamma - c\beta$
pairing).
\item \label{L2} \emph{Quillen-equivalence.} The pair
$(\Omegach_X,\Bbarch_X)$ is a Quillen equivalence between the
Francis--Gaitsgory--Rozenblyum model category on
$\Alg^{\fact,\aug,\comp}_X$ and the dual on
$\mathrm{CoAlg}^{\fact,\conil,\comp}_X$ (\cite[Chapter IV.5]{GR17}),
with unit and counit WEAK equivalences (in the model-category
sense, not necessarily chain-level qis). On $\Kosz(X)$ (H1)+(H2)+(H3),
weak equivalence promotes to chain-level qi via the
seven-fold TFAE hub (\texttt{ftm\_seven\_fold\_tfae\_platonic.tex}).
\item \label{L3} \emph{$(\infty,1)$-adjoint equivalence.} Passing to
$(\infty,1)$-localizations, the Quillen pair becomes an adjoint
equivalence of presentable stable $(\infty,1)$-categories. This is
the lane of Lurie-HA~\S5.5 factorization algebras; it works on the
conilpotent-complete locus UNCONDITIONALLY on (H3), because
$(\infty,1)$-localisation erases the distinction between
coderived and chain-level.
\item \label{L4} \emph{$(\infty,2)$-equivalence at properad level.} The
Francis--Gaitsgory properad-level $(\infty,2)$-adjoint equivalence
(Theorem~$A^{\infty,2}$, \texttt{theorem\_A\_infinity\_2.tex}
Theorem~\ref{thm:A-infinity-2}) lifts \ref{L3} along the
Hackney--Robertson inclusion $\Fact\mathrm{Op}\hookrightarrow\Fact\mathrm{Prop}$.
\end{enumerate}

\textbf{Resolution.} The statement ``$K^{2}\simeq\id$ on $\Kosz(X)$''
is UNAMBIGUOUSLY lane~\ref{L1} only on the standard archetypes
$\mathsf{G}/\mathsf{L}/\mathsf{C}/\mathsf{M}$ where (H3) is
witnessed (Lemma~\ref{lem:archetype-H123-witness} of the chapter).
Elsewhere it is~\ref{L2}--\ref{L4} with increasingly weak ``equals''.

\emph{Writing convention for the chapter.} Clause \ref{KR-ii} should
be read as \ref{L1}, witnessed archetype-by-archetype. Clause
\ref{KR-iii} (off $\Kosz(X)$) should be read as \ref{L3}, because the
Positselski coderived--contraderived equivalence is
$(\infty,1)$-categorical. The properad-level $(\infty,2)$ content is
Theorem~$A^{\infty,2}$ (lane \ref{L4}). The Quillen-equivalence
lane~\ref{L2} is load-bearing only at the promotion step between
\ref{KR-iii} and \ref{KR-ii} on the Koszul locus.

\emph{Load-bearing corollary.} The four archetype chain homotopies
$h_\mathsf{G}, h_\mathsf{L}, h_\mathsf{C}, h_\mathsf{M}$ are what
concretises $K^{2}\simeq\id$. They are not confabulated. They are:
\begin{itemize}
\item $h_\mathsf{G}$ (Heisenberg, $\cA=\mathcal{H}_k$):
the Koszul contraction on $\Bbarch(\mathcal{H}_k)=\Tc(s^{-1}\bar\mathcal{H}_k)$
is the standard simplicial contraction of the augmented tensor
coalgebra; $h_\mathsf{G}(\alpha_{-n_1}\cdots\alpha_{-n_p})
=\sum_j\alpha_{-n_1}\cdots\widehat{\alpha_{-n_j}}\cdots\alpha_{-n_p}\otimes\alpha_{-n_j}$
(up to the sign of the $s^{-1}$-desuspension), with
$[d_{\mathrm{bar}},h_\mathsf{G}]=\id$ on $\bar B^{\mathrm{ch}}_{\geq 1}(\mathcal{H}_k)$
and $h_\mathsf{G}|_{\bar B^{\mathrm{ch}}_0}=0$.
\item $h_\mathsf{L}$ ($V_k(\mathfrak{sl}_2)$ at $k\neq-2$):
the Sugawara-twisted contraction
$h_\mathsf{L}=h_\mathsf{G}+\tfrac{1}{2(k+h^\vee)}\sum_a[J^a,h_\mathsf{G}]J^a$,
a perturbation of the Heisenberg contraction by the Sugawara energy
tensor; the $[d,h_\mathsf{L}]$-identity is
$\id - p_\mathsf{L}$ with $p_\mathsf{L}$ the projection onto the
Fock-space vacuum modulo Sugawara descent.
\item $h_\mathsf{C}$ ($\beta\gamma_\lambda$):
the free-field OPE contraction $h_\mathsf{C}=\sum_n\beta_{-n}\partial_{\beta_{-n}}
+\gamma_{-n+\lambda-1}\partial_{\gamma_{-n+\lambda-1}}$ on the Fock space,
combined with $h_\mathsf{G}$ on the conformal-weight layers; nilpotency
of $\beta\partial_\beta+\gamma\partial_\gamma$ on the free-field vacuum
yields $[d,h_\mathsf{C}]=\id-p_\mathsf{C}$.
\item $h_\mathsf{M}$ ($\mathrm{Vir}_c$, weight-completed):
the Felder BRST contraction
$h_\mathsf{M}=Q^{-1}_{\mathrm{Felder}}\circ\pi^{\mathrm{ord}}$
on the Felder resolution of $\mathrm{Vir}_c$ at rational minimal-model
central charges, extended to generic $c$ via Fock-space continuation;
$[d,h_\mathsf{M}]=\id-p_\mathsf{M}$ with $p_\mathsf{M}$ the
projection onto the irreducible Virasoro module
$L(c,0)$, valid on the weight-completed ambient $\widehat{\mathrm{Vir}}_c$
per Remark~\ref{rem:class-M-pro-conilpotent}.
\end{itemize}

Each chain homotopy is independently checkable from the OPE of the
archetype algebra; none of them uses the wave-14 chapter's proof
machinery. They are the lane-\ref{L1} witnesses of $K^{2}\simeq\id$.

\subsection*{Cycle 2 (BEILINSON). Does the properad-level
$(\infty,2)$-arrow actually entail the chain-level Quillen equivalence?}

\textbf{Attack.} The chapter claims that Theorem~$A^{\infty,2}$
(Corollary~\ref{cor:eight-cor-properad-equiv}, line 722--744) at
properad level implies Theorem~A at chain level (Corollary~\ref{cor:classical-A-from-A-infinity-2},
line 1383--1409). This implication is by pullback along
$\{\mathrm{pt}\}\hookrightarrow\Ran(X)$ plus cohomological truncation
$\Fact(X)\to\Dmod(X)$. But: pullback along $\{\mathrm{pt}\}$ trivialises
the factorization structure, and cohomological truncation in general
does NOT preserve equivalences (the $(\infty,1)$-adjoint
equivalence at properad level may have nontrivial higher homotopies
whose cohomological truncation fails). Is the implication sound?

\textbf{Heal (DRINFELD).} Two correctives.

\emph{Corrective 1: pullback along $\{\mathrm{pt}\}$ is exact on the
conilpotent-complete locus.} On $\Kosz(X)$ with (H1)+(H2)+(H3), the
factorization algebra $\cA$ restricted to a single marked point
is the underlying vertex operator algebra, and the factorization
structure is recovered from the OPE (\cite{FBZ04}). The restriction
functor $\{\mathrm{pt}\}^*\colon\Fact(X)\to\Dmod(\{\mathrm{pt}\})$ is
exact on the (H1)+(H2)+(H3) locus because it is adjoint to the
constant factorization functor and both preserve colimits in
$\Fact(X)$ (\cite[\S4.3]{Francis2012}). So pullback sends the
$(\infty,2)$-equivalence at properad level to a chain-level
equivalence at the point, verbatim.

\emph{Corrective 2: cohomological truncation of an $(\infty,2)$-equivalence
yields an ordinary adjoint equivalence at the cohomological level.}
By the Bousfield--Friedlander theorem for $(\infty,2)$-categories,
truncation $\tau_{\leq 0}\colon\mathrm{Cat}_{(\infty,2)}\to\mathrm{Cat}_{(1,1)}$
preserves adjunctions and equivalences. So the $(\infty,2)$-adjoint
equivalence at properad level truncates to an ordinary
adjoint equivalence at chain-level (after H0). This is
what Corollary~\ref{cor:classical-A-from-A-infinity-2} claims.

\emph{Sharpened statement.} The implication
\begin{quote}
$(\infty,2)$-equivalence at properad level on $\Fact^{\aug}(X)$
$\;\Longrightarrow\;$ chain-level Quillen equivalence on
$\Alg^{\fact,\aug,\comp}_X$ at a marked point
\end{quote}
is a TWO-STEP implication:
\[
\underbrace{(\infty,2)\text{-equivalence}}_{\text{line~722--744}}
\;\xrightarrow[\tau_{\leq 0}]{\text{truncation}}\;
\underbrace{(\infty,1)\text{-equivalence}}_{\text{Lurie-HA~\S5.5, (L3)}}
\;\xrightarrow[\{\mathrm{pt}\}^*]{\text{point-restriction}}\;
\underbrace{\text{chain-level Quillen equivalence}}_{\text{(L1), Val16 applied}}.
\]
Each arrow is sound; the implication is not automatic but is
supplied by the cited machinery.

\emph{What MUST be added to the chapter.} Corollary~\ref{cor:classical-A-from-A-infinity-2}
should cite these two steps explicitly; its one-line proof at
line 1396--1409 currently invokes only the ``cohomological~$H^0$''
step, which is the second arrow but not the first. The first arrow
(the $(\infty,2)\to(\infty,1)$ truncation) should be added as a
separate line citing Bousfield--Friedlander or (more
canonically) \cite[Chapter I.3]{GR17} Proposition on
$(\infty,2)$-truncation.

\subsection*{Cycle 3 (BEILINSON). What is $\Kosz(C)$ precisely?
Which chiral algebras are chirally Koszul and which aren't?}

\textbf{Attack.} Theorem~\ref{thm:koszul-reflection} quantifies over
$\Kosz(X)\subset\Alg^{\fact,\aug,\comp}_X$. What is this locus,
PRECISELY?

\textbf{Heal (DRINFELD).} Defined at
\texttt{chapters/theory/chiral\_koszul\_pairs.tex:957--1016},
Definition~\ref{def:chiral-koszul-pair}. A chiral Koszul PAIR is
a pair $(\cA,\cC,\tau,F_\bullet)$ of:
\begin{itemize}
\item an augmented chiral algebra $\cA$ with chiral twisting datum
  $\tau\colon\bar B^{\mathrm{ch}}_X(\cA)\to\cA$ (Maurer--Cartan);
\item a conilpotent-complete dual coalgebra $\cC$ paired with $\cA$
  via Verdier duality $\mathbb{D}_{\Ran}(\cC)\simeq\Omegach_X(\cC^!)$;
\item a PBW filtration $F_\bullet\cA$ exhaustive, complete,
  bounded below, with associated graded $\mathrm{gr}^F\cA$ the
  natural polynomial chiral algebra on the generators.
\end{itemize}
$\Kosz(X)$ is then the image of chiral Koszul pairs in
$\Alg^{\fact,\aug,\comp}_X$ under first projection.

\emph{Members of $\Kosz(X)$.} The standard landscape (Heisenberg
$\mathcal{H}_k$, affine Kac--Moody $V_k(\mathfrak{g})$ at
$k\neq-h^\vee$, $\beta\gamma_\lambda$, Virasoro
$\mathrm{Vir}_c$ at generic $c$ (NOT at rational minimal-model
central charges without weight-completion), principal $\mathcal{W}_N$
at generic level, lattice VOAs of self-dual even lattice, Mukai-doubled
K3, free fermions). All satisfy the three recognition criteria:
\begin{enumerate}[label=\textup{(K\arabic*)}]
\item $\cA$ has PBW presentation with polynomial associated graded,
\item the associated graded is Koszul classical quadratic (in the
  sense of Priddy 1970 + Beilinson--Ginzburg--Soergel),
\item deformation (PBW) from $\mathrm{gr}^F\cA$ to $\cA$ is
  quasi-isomorphic to the chiral analogue in $\Fact(X)$ of the
  Loday--Vallette transfer.
\end{enumerate}

\emph{NON-members of $\Kosz(X)$.} (a) Vertex operator algebras with
non-quadratic relations obstructing PBW presentation (these are
the ``non-quadratic anomaly'' examples of
\texttt{non\_principal\_w\_duality.tex}); (b) chiral algebras at
fractional level admitting nontrivial spectral flow
(logarithmic CFT examples, fractional level $W_N$, non-admissible
representations of $\widehat{\mathfrak{g}}$); (c) certain W-algebras
constructed via non-principal Drinfeld--Sokolov reduction whose
Koszulness is established only in type A and conjectural elsewhere;
(d) any chiral algebra with MC4 obstruction (failure of
Maurer--Cartan to close at depth 4+) -- this is the genuinely
open frontier.

\emph{Non-circularity check.} $\Kosz(X)$ is defined by RECOGNITION
CRITERIA (K1)+(K2)+(K3), all checkable without invoking
Theorem~A. The theorem statement ``$K^{2}\simeq\id$ on $\Kosz(X)$''
is non-circular because $\Kosz(X)$ is defined antecedently, and
membership in $\Kosz(X)$ is verified on each archetype independently
in Lemma~\ref{lem:archetype-H123-witness}. The reverse direction
(every $\cA$ with $K^{2}\simeq\id$ is in $\Kosz(X)$) is the
seven-fold TFAE at \texttt{ftm\_seven\_fold\_tfae\_platonic.tex}.

\subsection*{Cycle 4 (BEILINSON). Is the symmetric-monoidal structure
compatible with factorisation (BD)?}

\textbf{Attack.} The statement ``symmetric-monoidal adjoint
equivalence'' at line 175 claims that $K$ commutes with some
symmetric-monoidal structure on both sides. On the algebra side this
is the $\star$-product on $\Fact(X)$ (Proposition~\ref{prop:fg-ambient-properties}).
On the coalgebra side it must be its dual. The Beilinson--Drinfeld
chiral pseudo-tensor $\otimes^{\ch}$ is NOT a bifunctor
(Remark~\ref{rem:why-fact}, line 832--847). Can the two tensor
structures be consistently paired?

\textbf{Heal (DRINFELD).} Three observations.

\emph{(a) The ambient for Theorem~A is not $(\Dmod(X),\otimes^{\ch})$,
it is $(\Fact(X),\star)$.} This is spelled out in
Definition~\ref{def:fg-ambient} (line 813--829) and Proposition~\ref{prop:fg-ambient-properties}.
$\star$ is a genuine bifunctor with strict unit, symmetric-monoidal
at $(\infty,2)$-level.

\emph{(b) The coalgebra side inherits $\star$ by duality.} The Verdier
duality $\mathbb{D}_{\Ran}$ is a $\star$-self-adjoint equivalence on
$\Fact(X)$ (\cite[IV.5, \S5.2]{GR17}). The pair
$(\mathrm{CoFact}^{\conil,\comp}_X,\star^\vee)$ with
$\star^\vee:=(\star)^{\mathrm{op}}\circ\mathbb{D}_{\Ran}^{\otimes 2}$
is symmetric-monoidal at $(\infty,2)$-level, and
$\Bbarch_X$ / $\Omegach_X$ both preserve the two tensor products up to
canonical natural isomorphism. This is the content of
\cite[\S4.2]{Francis2012} combined with the chiral Verdier duality
of Theorem~\ref{thm:bar-cobar-verdier}.

\emph{(c) Compatibility with Beilinson--Drinfeld factorisation
structure.} On $\Fact(X)$, factorisation structure is BUILT INTO
the objects: a factorization algebra is a $\star$-algebra by
definition. The Koszul reflection $K$ sends a factorization algebra
to a factorization coalgebra (not a bimodule, not a
``two-sided'' factorization object): the factorisation compatibility
is the statement that, on disjoint point-configurations $z_1,\ldots,z_n$,
the bar complex $\Bbarch_X(\cA)(z_1,\ldots,z_n)$ factorises as
$\bigotimes_i\Bbarch_X(\cA)(z_i)$ (the factorisation being the
$\star$-tensor of the pointwise bar complexes). Dually for
$\Omegach_X$.

\emph{Sharpened statement.} ``Symmetric-monoidal adjoint equivalence''
in clause \ref{KR-i} should be qualified: it is SM at the
$(\infty,2)$-categorical level on $(\Fact(X),\star)$, which is the
ambient in which the statement makes sense. The chain-level
specialisation at a marked point recovers the classical
symmetric-monoidal adjoint equivalence of Loday--Vallette (line
1004--1017, \ref{A2-ii}).

\emph{What MUST be added.} Theorem~\ref{thm:koszul-reflection} should
say explicitly ``symmetric-monoidal at $(\infty,2)$-level on
$(\Fact(X),\star)$'' rather than an unqualified ``symmetric-monoidal''.
The cross-volume factorization-compatibility is load-bearing
(downstream: Vol III $\Phi_d$ is SM, \cite{cy_to_chiral.tex}; the
Theorem~A SM structure is what makes $\Phi_d$-compatibility
coherent).

\subsection*{Cycle 5 (BEILINSON). Does Francis--Gaitsgory 2012 directly
give $K^{2}\simeq\id$, or does it need supplementing?}

\textbf{Attack.} The chapter cites Francis--Gaitsgory~\cite{Francis2012}
for the $\star$-product adjunction and GR17~\cite{GR17} for the
$(\infty,2)$-enhancement. Does either paper establish
$K^{2}\simeq\id$ directly?

\textbf{Heal (DRINFELD).} Neither. Here is what each supplies.

\emph{Francis 2012 supplies.} The $\star$-monoidal structure on
$\Fact(X)$ (\cite[\S4.2]{Francis2012}), the universal-property
characterisation of factorization algebras as $\star$-algebras
(\cite[\S3]{Francis2012}), and the chiral Deligne conjecture
(\cite[\S5]{Francis2012}) which gives the brace algebra
structure on $\RHom_{\cA\otimes\cA^{\mathrm{op}}}(\cA,\cA)$.

\emph{GR17 supplies.} The $(\infty,2)$-categorical sheaf-of-categories
formalism (\cite[Part IV]{GR17}), the $(\infty,2)$-enhancement of
$\Fact(X)$ (\cite[Chapter IV.5, Theorem 3.1.2]{GR17}), and the
correspondence / base-change formalism (\cite[Chapter IV.1]{GR17}).

\emph{Neither paper supplies $K^{2}\simeq\id$ directly.} The Koszul
reflection is stated in the wave-14 chapter as a THEOREM, proved by
combining the above machinery with:
\begin{itemize}
\item \textbf{Vallette 2016.} Quillen equivalence between
augmented $A_\infty$-algebras and conilpotent CDG-coalgebras
(\cite[Theorem 2.1]{Val16}) -- this is the Quillen lane.
\item \textbf{Positselski 2011+2018.} Coderived $\simeq$
contraderived equivalence (\cite{Positselski2011}) and
weakly-curved $A_\infty$ $\simeq$ conilpotent CDG
(\cite[Theorem 9.1]{Positselski2018}) -- this is the
curvature-aware lane.
\item \textbf{Hinich 2003.} Symmetric-monoidal Quillen equivalence
(\cite[\S6]{Hinich2003}) -- this is the SM promotion.
\item \textbf{Loday--Vallette 2012.} Chain-level bar-cobar at a
marked point (\cite[Theorem 11.4.1]{LV12}) -- this is the
chain-level verification at $\{\mathrm{pt}\}$.
\item \textbf{Hackney--Robertson 2019.} $\infty$-properad model
structure (\cite[Theorem 5.12]{HackneyRobertson2019}) -- this
is the properad lift.
\end{itemize}

\emph{Sharpened statement.} $K^{2}\simeq\id$ on $\Kosz(X)$ is a
THEOREM of the present programme, not a corollary of a single cited
source. Its proof is a synthesis of the above five machinery
packages, with the chiral specialisation supplied by
Beilinson--Drinfeld~\cite[\S3.6]{BD04} (chiral operad Koszulness).

\emph{What the chapter correctly attributes.} The proof of
Theorem~\ref{thm:koszul-reflection} at line 203--244 lists Vallette,
Positselski, Hinich, and Francis--Gaitsgory correctly; no source
is mis-attributed. The INDEPENDENT VERIFICATION sources at
Remark~\ref{rem:ainf2-hziv} (Francis chiral Deligne, Lurie HA~\S5.5,
Hackney--Robertson) are genuinely disjoint from the derivation
sources, hence legitimate independent paths.

\subsection*{Cycle 6 (BEILINSON). Does the consolidation lose content
vs the original eight subclaims?}

\textbf{Attack.} The chapter advertises ``eight prior parallel
statements $\to$ one theorem'' (line 54--63). If the consolidation
is lossy -- if any of the eight subclaims carries content not
recoverable from $K^{2}\simeq\id$ -- then the consolidation is
false simplification.

\textbf{Heal (DRINFELD).} The eight subclaims, as recovered in
Corollaries~\ref{cor:eight-cor-twisting-morphism}--\ref{cor:eight-cor-properad-equiv}
(line 546--744), are:
\begin{enumerate}
\item Twisting morphism $\tau$ is a Maurer--Cartan element
(\texttt{cor:eight-cor-twisting-morphism}, line 546--571).
\item Counit qi on $\Kosz(X)$ (\texttt{cor:eight-cor-counit-qi},
line 573--587).
\item Unit weak equivalence on $\Conil(X)$ (\texttt{cor:eight-cor-unit-weq},
line 589--603).
\item Twisted tensor acyclicity on $\Kosz(X)$
(\texttt{cor:eight-cor-twisted-tensor}, line 605--629).
\item Bar concentration in weight 1 at $g=0$ on class $\mathsf{G}$
(\texttt{cor:eight-cor-bar-weight-1}, line 631--661).
\item SC-formality on class $\mathsf{G}$ only
(\texttt{cor:eight-cor-sc-formality}, line 663--688).
\item $R$-twisted $\Sigma_n$-descent (\texttt{cor:eight-cor-R-descent},
line 690--720).
\item Francis--Gaitsgory $(\infty,2)$-equivalence at properad level
(\texttt{cor:eight-cor-properad-equiv}, line 722--744).
\end{enumerate}

\emph{Each is recovered verbatim or in sharpened scope.} Let us check
one by one:
\begin{itemize}
\item (1) is a restatement of clause \ref{KR-i} at the chain level on
a specific element (the twisting morphism as MC).
\item (2)-(3) are the two halves of clause \ref{KR-ii} on $\Kosz(X)$.
\item (4) is a direct consequence of (2): twisted tensor acyclicity
is equivalent to the counit being a qi (LV12 Theorem 2.4.1).
\item (5)-(6) are class-$\mathsf{G}$ specialisations via the seven-fold
TFAE + the $\kappa=0$ condition characterising class $\mathsf{G}$.
These are STRICTLY WEAKER than the general $K^{2}\simeq\id$
statement (they hold on class $\mathsf{G}$ only), so the
consolidation into a single statement with class-$\mathsf{G}$ as
a specialisation is lossless.
\item (7) is the ordered-bar / symmetric-bar comparison, stated at
the level of the $L_R$-twisted descent, with unitarity
hypothesis $R(z)R^{\mathrm{op}}(-z)=\id$.
\item (8) is the properad-level $(\infty,2)$-enhancement at
Theorem~$A^{\infty,2}$ (line 955--1040).
\end{itemize}

\emph{Is any content lost?} CHECK: the original statement of
``bar concentration in weight 1'' at $g\geq 1$ on class $\mathsf{G}$
was in pre-wave-14 Vol I. Is it preserved? YES, in
Corollary~\ref{cor:eight-cor-bar-weight-1} line 644--647: at $g\geq 1$,
the fiberwise curvature $d^2_{\mathrm{fib}}=\kappa\cdot\omega_g$
obstructs concentration in weight 1 UNLESS $\kappa=0$, which holds in
class $\mathsf{G}$. So the extension to $g\geq 1$ is
class-$\mathsf{G}$-only, matching the original statement.

CHECK: the original ordered bar $R$-twisted descent was proved under
unitarity $R(z)R^{\mathrm{op}}(-z)=\id$. Is it preserved? YES,
Lemma~\ref{lem:R-twisted-descent} line 1421--1519 retains unitarity
as hypothesis (R2) with explicit scope at
Remark~\ref{rem:R-descent-unitarity-scope} (rational Yangians,
trigonometric $U_q(\hat\fg)$ at generic $q$, elliptic Belavin/Felder
scope-dependent).

CHECK: the original Francis--Gaitsgory $(\infty,2)$-equivalence was
claimed at properad level. Is it preserved? YES, Theorem~\ref{thm:A-infinity-2}
line 955--1040 with the four-step proof.

\emph{Conclusion.} The consolidation into a single theorem is
LOSSLESS. The eight subclaims are all recoverable as corollaries;
no content is thrown away. The unification is a GAIN in
clarity (one theorem, one invariant $K$, one slogan
$K^{2}\simeq\id$) without loss of specificity (each corollary states
its specialisation cleanly).

\subsection*{Cycle 7 (BEILINSON). The MC4 regime: exhibit a specific
$\cA$ where (H3) fails and $K^{2}\neq\id$ at chain level.}

\textbf{Attack.} The chapter claims at line 95--97 that MC4 is the
``only genuinely open frontier'', and at line 1970--1984 that
$\Pi 4$ (unbounded-rank Koszul reflection) is the corresponding
conjecture. Can a specific $\cA$ be exhibited with (H3) failing and
$K^{2}\neq\id$ at chain level?

\textbf{Heal (DRINFELD).} Yes. The specific example is
\emph{critical-level affine Kac--Moody}: $\cA=V_{-h^\vee}(\mathfrak{g})$
with $\mathfrak{g}$ simple, at the critical level $k=-h^\vee$.

\emph{Why (H3) fails.} At critical level, the Feigin--Frenkel center
$\mathfrak{z}(\widehat{\mathfrak{g}})=\mathrm{Cent}(V_{-h^\vee}(\mathfrak{g}))$
is (classically) an infinite-dimensional polynomial algebra
generated by the higher Casimirs $S_2,S_3,\ldots,S_{h^\vee}$ (Sugawara,
higher Sugawara, $\mathcal{W}$-algebra generators). The bar complex
$\Bbarch_X(V_{-h^\vee}(\mathfrak{g}))$ inherits this center as a
subcomplex, and the conformal weight grading on the center extends
to unbounded depth because the higher Casimirs have unbounded
conformal weight with unbounded multiplicity (the polynomial ring
$\Bbarch_X(\mathrm{Sym}(\mathfrak{z}))$ has finite-dim conformal-weight
pieces at each weight BUT the bar direction opens up unbounded
weight differentials via the Sugawara descent).

More precisely: at $k=-h^\vee$, the Sugawara tensor
$L_n^{\mathrm{Sugawara}}=\tfrac{1}{2(k+h^\vee)}\sum_a:J^a J^a:_n$ has
a POLE at $k=-h^\vee$; the Sugawara vector becomes central
($=\mathrm{Cent}$), and the Virasoro action degenerates. The bar
complex $\Bbarch_X(V_{-h^\vee})$ now carries bar differentials whose
higher Massey products $m_k$ do not stabilise at any finite $k$
(they are the higher Sugawara descendants, unbounded in $k$); (H3)
fails because the conformal-weight pieces of the bar complex at
weight $w$, although individually finite, do not satisfy the
finiteness condition AFTER Sugawara descent because the Sugawara
tower connects them with unbounded depth.

\emph{What fails and what survives.}
\begin{itemize}
\item The chain-level qi $\varepsilon_\cA\colon\Omegach_X\Bbarch_X(\cA)\to\cA$
DEGRADES at critical level: the counit is no longer a qi at finite
conformal weight, only at the pro-conilpotent completion of
$V_{-h^\vee}$ (the Feigin--Frenkel / Kac-Kazhdan completion). This is
the standard story: Feigin--Frenkel \cite{FF92} show that the
naive direct-sum completion is $V_{-h^\vee}$ itself, but the
Koszul-reflective completion is a proper enlargement, the
``completed $V_{-h^\vee}$'' whose derived center IS the classical
Langlands-dual Lie algebra $\mathfrak{z}(\widehat{\mathfrak{g}})$.
\item The $(\infty,1)$-equivalence (lane \ref{L3}) SURVIVES: the
coderived-contraderived equivalence of Positselski \cite{Positselski2011}
holds without (H3), because $(\infty,1)$-localisation absorbs the
Mittag-Leffler failure. This is what clause \ref{KR-iii} claims.
\item The $(\infty,2)$-equivalence at properad level (lane \ref{L4})
SURVIVES on the conilpotent-complete locus WITHOUT (H3): Theorem~$A^{\infty,2}$
does not require (H3), only (H1)+(H2). The conilpotent-complete
locus IS the honest ambient for critical-level affine KM after
Feigin--Frenkel completion.
\end{itemize}

\emph{Sharpened conclusion.} At critical level $k=-h^\vee$:
\begin{itemize}
\item lane \ref{L1} (chain-level): $K^{2}\simeq\id$ FAILS (because
(H3) fails).
\item lane \ref{L2} (Quillen): $K^{2}\simeq\id$ HOLDS at Quillen-equivalence
level after pro-conilpotent completion (Feigin--Frenkel).
\item lane \ref{L3} ($(\infty,1)$): $K^{2}\simeq\id$ HOLDS in the
coderived category unconditionally.
\item lane \ref{L4} ($(\infty,2)$-properad): $K^{2}\simeq\id$ HOLDS
unconditionally on conilpotent-complete factorization properads.
\end{itemize}

So the MC4-regime failure is genuine at chain level but does NOT
propagate upward: the weaker equivalences (lanes \ref{L2}--\ref{L4})
are preserved. The honest claim ``$K^{2}\simeq\id$'' at critical
level is scoped to the coderived / $(\infty,1)$-categorical lane,
which is exactly what clause \ref{KR-iii} of the theorem captures.
The chapter is correct.

\emph{What MUST be added.} An example remark immediately after
Remark~\ref{rem:class-M-pro-conilpotent} (line 477--505) spelling
out critical-level affine KM as the canonical (H3)-failure witness.
This complements the existing $\Pi 4$ conjecture at line 1970--1984
by EXHIBITING the obstruction on a specific example, not just
naming the conjecture.

%================================================================
\section*{Distilled heal: step-by-step first-principles proof of $K^{2}\simeq\id$ on $\Kosz(C)$}
%================================================================

Combining the seven cycle outcomes, the following is the rigorous
first-principles chain-level proof of $K^{2}\simeq\id$ on
$\Kosz(X)$, with all lane-discipline markers in place.

\begin{theorem}[Chain-level $K^{2}\simeq\id$ on $\Kosz(X)$;
healed statement]
\label{thm:healed-k-squared-id}
Let $X$ be a smooth projective curve over a field of characteristic 0.
Let $\cA\in\Kosz(X)\subset\Alg_X^{\fact,\aug,\comp}$
(Definition~\ref{def:chiral-koszul-pair}) satisfy (H1)+(H2)+(H3)
(\texttt{theorem\_A\_infinity\_2.tex:69--87}). Then:
\begin{enumerate}[label=\textup{(\roman*)}]
\item The counit $\varepsilon_\cA\colon\Omegach_X\Bbarch_X(\cA)
\simequiv\cA$ is a chain-level quasi-isomorphism, witnessed by a
named chain homotopy
$h_\cA\colon\Omegach_X\Bbarch_X(\cA)\to\Omegach_X\Bbarch_X(\cA)[-1]$
satisfying $[d,h_\cA]=\id-p_\cA$ where
$p_\cA=(\varepsilon_\cA)^{-1}\circ\varepsilon_\cA$ is the retraction.
\item The unit $\eta_\cA\colon\cA\simequiv\Bbarch_X\Omegach_X(\cA)$
is a chain-level quasi-isomorphism, witnessed by a named chain
homotopy $\tilde h_\cA$ satisfying $[d,\tilde h_\cA]=\id-\tilde p_\cA$.
\item The composite $K^{2}=(\Omegach_X\circ\Bbarch_X)\circ(\Omegach_X\circ\Bbarch_X)$
is chain-level quasi-isomorphic to $\id_{\Alg_X^{\fact,\aug,\comp}}$
restricted to $\Kosz(X)$.
\end{enumerate}
\end{theorem}

\begin{proof}[Step-by-step first-principles proof]

\emph{Step 1 (Chain complexes are well-defined on $\Kosz(X)$).}
By Definition~\ref{def:chiral-koszul-pair}, every $\cA\in\Kosz(X)$
admits a chiral twisting datum
$\tau_\cA\colon\bar B^{\mathrm{ch}}_X(\cA)\to\cA$ and a PBW filtration
$F_\bullet\cA$ exhaustive, complete, bounded below. The bar
$\Bbarch_X(\cA):=T^c_{(\Fact(X),\star)}(s^{-1}\bar\cA)$ with
alternating-sign deconcatenation differential is a
conilpotent factorization coalgebra (Step 1 of
Theorem~\ref{thm:A-infinity-2} proof at line 1049--1057). The cobar
$\Omegach_X(\cC):=T_{(\Fact(X),\star)}(s^{-1}\bar\cC)$ with twisted
differential $d_\cC+\partial_\tau$ is a factorization algebra (line
1059--1067). Both are chain complexes in $\Fact(X)$.

\emph{Step 2 ($(\infty,1)$-adjunction in $\Fact(X)$).}
Adjunction $\Bbarch_X\dashv\Omegach_X$ is the universal property of
free and cofree constructions in the $\star$-monoidal ambient
$(\Fact(X),\star)$ (Step 1 of Theorem~$A^{\infty,2}$ proof at line
1062--1067, citing \cite[IV.5, \S2.4]{GR17}). This is a
Quillen adjunction (lane \ref{L2}) promoting to an
$(\infty,1)$-adjunction (lane \ref{L3}).

\emph{Step 3 (Unit/counit on the conilpotent-complete locus).}
The unit $\eta_\cA\colon\cA\to\Omegach_X\Bbarch_X(\cA)$ is the
standard bar-cobar bicomplex contraction; acyclicity on
$\Alg^{\fact,\aug,\comp}_X$ follows from a spectral-sequence
argument (Step 2 of Theorem~$A^{\infty,2}$ proof at line 1071--1083).
The bar-length filtration on $\Bbarch_X(\cA)$ is exhaustive
by completeness (H2), and each associated graded piece
$\mathrm{gr}^p\Bbarch_X(\cA)$ carries the classical Koszul complex for
$(\Ass,\Ass^!)$ via cohomological truncation and
point-restriction (line 1075--1079); the Koszul complex is
acyclic over the unit by \cite[Theorem~7.4.1]{LV12}.
Exactness of the base-change propagates acyclicity to the
$E_1$-page.

Convergence of the spectral sequence to the unit morphism is
Mittag-Leffler for the bar-length filtration under
conilpotent-completeness (H2), supplied by the cofinal tower
$\{\Bbarch_X(\cA/\bar\cA^{\,n})\}_{n\geq 1}$ (Remark~\ref{rem:A-inf-2-ml-step-2}
at line 1135--1161). Hypothesis (H3) supplies surjectivity of the
stage maps in every cohomological degree; Mittag-Leffler then
yields chain-level convergence. Counit is dual.

\emph{Step 4 (Promotion to chain-level qi on $\Kosz(X)$).}
The seven-fold TFAE at $E_2$-collapse hub
(\texttt{ftm\_seven\_fold\_tfae\_platonic.tex}
Theorem~\ref{thm:ftm-seven-fold-tfae-via-hub-spoke}) states that on
genus-0 $X$ and $\cA\in\Kosz(X)$ satisfying PBW filtration
conditions, the following are equivalent:
\begin{enumerate}[label=\textup{(\roman*$^*$)}]
\item $\cA$ is chirally Koszul;
\item $\varepsilon_\cA$ is a chain-level qi (counit qi);
\item $\eta_\cA$ is a chain-level qi (unit weak equivalence);
\item the twisted tensor complex is acyclic;
\item the bar complex concentrates in weight 1 (on class $\mathsf{G}$
at all $g$; on other classes at $g=0$ only);
\item the PBW bar-cobar spectral sequence collapses at $E_2$.
\end{enumerate}
Under (H1)+(H2)+(H3), each of these holds; the counit qi at (ii$^*$)
and the unit qi at (iii$^*$) are the two halves of $K^{2}\simeq\id$
at chain level.

\emph{Step 5 (Archetype witnesses of the chain homotopy).}
For each archetype $\mathsf{G}/\mathsf{L}/\mathsf{C}/\mathsf{M}$,
the chain homotopy $h_\cA$ is constructed explicitly:
\begin{itemize}
\item \textbf{Heisenberg $\mathsf{G}$}: $\cA=\mathcal{H}_k$. The bar
complex $\bar B^{\mathrm{ch}}(\mathcal{H}_k)=T^c(s^{-1}\bar\mathcal{H}_k)$
is generated by $\alpha_{-n_1}\cdots\alpha_{-n_p}$ with
$n_j\geq 1$. The Koszul contraction
\[
h_\mathsf{G}(\alpha_{-n_1}\cdots\alpha_{-n_p})
=\sum_{j=1}^{p}(-1)^{j-1}\alpha_{-n_1}\cdots\widehat{\alpha_{-n_j}}
\cdots\alpha_{-n_p}\otimes\alpha_{-n_j}
\]
(with appropriate $s^{-1}$-sign conventions) satisfies
$[d_{\mathrm{bar}},h_\mathsf{G}]=\id$ on
$\Bbarch^{\geq 1}(\mathcal{H}_k)$ and
$h_\mathsf{G}|_{\Bbarch^{0}}=0$. Direct verification: $d_{\mathrm{bar}}$
is the alternating-sign cut-and-contract, $h_\mathsf{G}$ is the
alternating-sign insert-and-uncontract; the combinatorial identity
is the standard simplicial contraction of the augmented tensor
coalgebra (LV12 Proposition 2.1.12).
\item \textbf{Affine Kac--Moody $\mathsf{L}$}: $\cA=V_k(\mathfrak{sl}_2)$
at $k\neq-2$. The Sugawara-twisted chain homotopy is
\[
h_\mathsf{L}=h_\mathsf{G}+\frac{1}{2(k+h^\vee)}\sum_a[J^a,h_\mathsf{G}]\cdot J^a,
\]
with $h_\mathsf{G}$ the Heisenberg contraction on the
generating $\fg$-mode factor and the second term the Sugawara
descent contribution. $[d,h_\mathsf{L}]=\id-p_\mathsf{L}$ where
$p_\mathsf{L}$ is the projection onto the irreducible vacuum module
$L(k,0)$; verification uses the Gelfand--Fuks cohomology computation
for $\widehat{\mathfrak{sl}_2}$.
\item \textbf{$\beta\gamma$ $\mathsf{C}$}: $\cA=\beta\gamma_\lambda$.
The free-field OPE chain homotopy
\[
h_\mathsf{C}=\sum_{n\geq 1}\bigl(\beta_{-n}\cdot\partial_{\beta_{-n}}
+\gamma_{-n+\lambda-1}\cdot\partial_{\gamma_{-n+\lambda-1}}\bigr)
\]
on the $\beta\gamma$ Fock space, combined with
$h_\mathsf{G}$ on conformal-weight layers. $[d,h_\mathsf{C}]=\id-p_\mathsf{C}$
where $p_\mathsf{C}$ is the projection onto the free-field vacuum.
\item \textbf{Virasoro $\mathsf{M}$}: $\cA=\mathrm{Vir}_c$. The Felder
BRST chain homotopy
\[
h_\mathsf{M}=Q_{\mathrm{Felder}}^{-1}\circ\pi^{\mathrm{ord}}
\]
on the Felder resolution of $\mathrm{Vir}_c$ at rational minimal-model
central charges, extended to generic $c$ via Fock-space
continuation (working on the weight-completed
$\widehat{\mathrm{Vir}}_c=\varprojlim_N\mathrm{Vir}_c/\mathrm{Vir}_c^{\geq N}$,
Remark~\ref{rem:class-M-pro-conilpotent}). $[d,h_\mathsf{M}]=\id-p_\mathsf{M}$
with $p_\mathsf{M}$ the projection onto $L(c,0)$.
\end{itemize}

\emph{Step 6 (Conclusion).}
The four archetype chain homotopies $h_\mathsf{G},h_\mathsf{L},h_\mathsf{C},h_\mathsf{M}$
witness $K^{2}\simeq\id$ at chain level (lane \ref{L1}) on the
representative of each archetype. By
Lemma~\ref{lem:archetype-H123-witness}, each representative is in
$\Kosz(X)$ with (H1)+(H2)+(H3) satisfied on the
conformal-weight grading (possibly after pro-conilpotent
weight completion for class $\mathsf{M}$). The chain-level
$K^{2}\simeq\id$ therefore holds on the entire standard
landscape; extension to abstract $\cA\in\Kosz(X)$ is by the
seven-fold TFAE of Step 4 plus the Vallette/Positselski/Hinich
machinery.

Off $\Kosz(X)$, lanes \ref{L2}--\ref{L4} survive; lane \ref{L1}
degrades to the coderived statement, which is the content of
clause \ref{KR-iii}.
\end{proof}

%================================================================
\section*{Hidden structure exposed by the cycles}
%================================================================

The seven cycles, taken together, expose a structural feature of
the programme that was previously distributed across six files:

\begin{enumerate}[label=\textbf{(S\arabic*)}]
\item \textbf{Four-lane discipline is load-bearing.}
$K^{2}\simeq\id$ is not a single theorem; it is four theorems at
increasing categorical strength (lanes \ref{L1}--\ref{L4}), each
with its own scope and witnesses. The chapter should state which
lane each corollary (D1)--(D14) of
Corollary~\ref{cor:ainf2-downstream-list} lives in.
\item \textbf{The archetype chain homotopies are a genuine witness
table.} The four homotopies $h_\mathsf{G},h_\mathsf{L},h_\mathsf{C},h_\mathsf{M}$
are the concrete content of $K^{2}\simeq\id$ on the four classes.
They complement Lemma~\ref{lem:archetype-H123-witness} (which
witnesses (H1)+(H2)+(H3) on the bar COMPLEX side) by providing
the chain-level HOMOTOPY witnesses on the derived-identity side.
\item \textbf{Critical-level affine KM is the canonical
(H3)-failure.} The $\Pi 4$ conjecture on unbounded-rank Koszul
reflection (line 1970--1984) has a specific example: critical-level
$V_{-h^\vee}(\mathfrak{g})$. This example belongs after
Remark~\ref{rem:class-M-pro-conilpotent} or in a new remark
attached to $\Pi 4$.
\item \textbf{The pullback-along-point argument from
$(\infty,2)$ to chain-level is TWO arrows, not one.} The step from
properad $(\infty,2)$-equivalence to chain-level Quillen equivalence
factors through the $(\infty,2)\to(\infty,1)$-truncation, which is
cited (via Bousfield--Friedlander / GR17) but not named explicitly
in Corollary~\ref{cor:classical-A-from-A-infinity-2}. This is a
minor expository gap.
\item \textbf{Symmetric-monoidal compatibility needs the
$\star$-qualifier.} ``Symmetric-monoidal adjoint equivalence'' at
Theorem~\ref{thm:koszul-reflection} should read
``$(\Fact(X),\star)$-symmetric-monoidal adjoint equivalence''. This
is a scope marker, not a theorem change.
\item \textbf{The Vol I / Vol III bridge through $\Phi_d$ uses
the SM structure of $K$.} Vol III's two-stage factorisation
$\Phi_d=\mathrm{Sp}^{\mathrm{ch}}_{\Sigma_{d-1},C}\circ\Phi^{\mathrm{FA}}_d$
composes with the SM structure of $K$ on $\Fact(C)$. The
compatibility is the content of Corollary~(D9)
(Corollary~\ref{cor:ainf2-downstream-list} line 1630--1645); the
$(\infty,2)$-functoriality of $\Phi_d$ PASSES THROUGH
Theorem~$A^{\infty,2}$, making Vol III's universal trace identity
consistent with Vol I's conductor identity.
\end{enumerate}

%================================================================
\section*{Open frontiers}
%================================================================

Three genuine opens remain after the seven cycles.

\begin{enumerate}[label=\textbf{(O\arabic*)}]
\item \textbf{Modular-family extension over $\overline{\mathcal{M}}_{g,n}$}
(Remark~\ref{rem:A-infinity-2-modular-family-scope}, line 1163--1353).
The $[\alpha_{\mathrm{BC}}]$ base-change class and
$[\beta_{\mathrm{nodal}}]$ nodal-sewing class remain unverified.
Vol I Theorem~$A^{\infty,2}$ is proved on a fixed curve $X$;
extension to the universal stable family is cited only at
GR17 Part IV + Mok25 (2025 arXiv). Clutching-uniqueness for
Theorem~D at $g\geq 2$ depends on clearing $[\beta_{\mathrm{nodal}}]$;
currently conditional on Class $\mathsf{G}$ (unconditional) and
Class $\mathsf{L}$ (conditional on generic non-critical level);
Class $\mathsf{M}$ weight-completed inherits from Vol II
\texttt{universal\_celestial\_holography.tex}.
\item \textbf{Lagrangian--Koszul converse $\Pi 3$.} Lagrangian
transversality of $M_{\mathrm{comp}}$ at $[\cA]$ implies chiral
Koszulness. Stated as Conjecture~\ref{conj:lagrangian-koszul-converse}
(line 1962--1968); requires the chiral PTVV framework which is in
Vol III. The converse (chiral Koszulness $\Rightarrow$ Lagrangian
transversality) is proved; the forward is open.
\item \textbf{Unbounded-rank Koszul reflection $\Pi 4$.} On
non-(H3) loci (e.g., critical-level affine KM after pro-conilpotent
completion where the completion has unbounded-depth Sugawara
tower), the chain-level $K^{2}\simeq\id$ fails but the coderived
holds. A Mittag-Leffler-type condition strengthening (H3) but
weaker than finite-dim is conjecturally the correct bridging
hypothesis. Stated as Conjecture~\ref{conj:unbounded-rank-koszul-reflection}
(line 1970--1984).
\end{enumerate}

%================================================================
\section*{Cross-volume harmonies (Drinfeld unifying principle)}
%================================================================

The healed $K^{2}\simeq\id$ on $\Kosz(X)$ couples to the three
volumes as follows.

\begin{itemize}
\item \textbf{Vol I (algebra).} Theorem~A is the backbone adjunction
$\Omegach_X\dashv\Bbarch_X$. Its healed statement (four lanes
\ref{L1}--\ref{L4}) gives the four operational levels at which the
programme's downstream theorems (B Positselski, C derived-centre
complementarity, D obstruction-tower universality, H Hochschild
concentration) can be inscribed. In particular, Theorem~D's
$\mathrm{obs}_g=\kappa\cdot\lambda_g$ is currently inscribed at
chain level (lane \ref{L1}) on Class $\mathsf{G}$ and conditionally
on Class $\mathsf{L}$ via the clutching-uniqueness logic which
depends on $[\beta_{\mathrm{nodal}}]$.
\item \textbf{Vol II (physics).} Universal Holography
$\Phi_{\mathrm{hol}}\colon\mathrm{ChirAlg}^{\omega,BL}_X\to\mathrm{HT\text{-}QFT}_{X\times\R}$
at \texttt{universal\_holography.tex} uses the chain-level
$K^{2}\simeq\id$ to identify bulk = $Z^{\mathrm{der}}_{\mathrm{ch}}(\cA)$
with boundary = $\cA$ via the brace / $\mathrm{SC}^{\mathrm{ch,top}}$
action. The bulk-boundary consistency IS $K^{2}\simeq\id$
evaluated on the derived-centre pair. Topologisation of
$\mathrm{SC}^{\mathrm{ch,top}}$ to $E_3$-topological (at
noncritical level for Class $\mathsf{L}$, weight-completed for
Class $\mathsf{M}$) is THE chiral shadow of Lurie's
HA~\S5.5 result that $E_3$-algebras = braided monoidal categories
with Drinfeld centre.
\item \textbf{Vol III (geometry).} The two-stage factorisation
$\Phi_d=\mathrm{Sp}^{\mathrm{ch}}_{\Sigma_{d-1},C}\circ\Phi^{\mathrm{FA}}_d$
at \texttt{cy\_to\_chiral.tex} produces an $E_1$-chiral algebra on
the reference curve $C$ from a CY$_d$ category. Theorem~A
($K^{2}\simeq\id$) is the backbone of the bar--cobar duality on
that curve; a SINGLE CY$_d$ category yields a FAMILY of chiral
shadows indexed by $(\Sigma_{d-1},C)$, each satisfying
$K^{2}\simeq\id$. The universal trace identity
$K\circ\Phi=\kappa_{\mathrm{ch}}$ is the commutativity of the
single square coupling the Vol I conductor to the Vol III
Borcherds weight $\kappa_{\mathrm{BKM}}$ on K3-fibered Class A.
\end{itemize}

The inner music: $K$ is the invariant; $\Bbarch_X$ and
$\Omegach_X$ are coordinate functorial expressions of it.
$K^{2}\simeq\id$ is Pontryagin-style involutivity of this
invariant, and the four-lane discipline exhibits four
progressively weaker notions of equality, each appropriate
for a different stratum of the Koszul locus.

One symmetry, one diagram, one equation: $K^{2}\simeq\id$.

\end{document}
